
./calculusApplicationsOfDerivatives.tex:7544:\حصہ{خط بندی اور تفرقات}\شناخت{حصہ_استعمال_خط_بندی_اور_تفرقات}
./calculusApplicationsOfDerivatives.tex:7766:\جزوحصہء{تفرقات}
./calculusApplicationsOfDerivatives.tex:7789:حاصل کرتے ہیں۔ یہ مساوات کہتی ہے کہ \عددی{\dif x\ne 0} کی صورت میں \عددی{f'(x)} \اصطلاح{تفرقات }\فرہنگ{تفرقات}\حاشیہب{differentials}\فرہنگ{differentials} \عددی{\dif y} اور \عددی{\dif x}  کا حاصل تقسیم ہو گا۔
./calculusApplicationsOfDerivatives.tex:7827:\جزوحصہء{تفرقات کی مدد سے تبدیلی کی اندازاً قیمت}
./calculusApplicationsOfDerivatives.tex:7946:\عددی{r} اور \عددی{H} کے تفرقات کا تعلق لکھتے ہیں۔
./calculusApplicationsOfDerivatives.tex:8279:\موٹا{تفرقات}\\
./calculusApplicationsOfIntegrals.tex:4710:یہ دکھانے کی خاطر ہم مساوات \حوالہ{مساوات_تکمل_استعمال_سطحی_رقبہ_الف} کو وقفہ \عددی{[a,b]} پر کسی تفاعل  کا ریمان مجموعہ لکھتے ہیں۔لمبائی قوس کے حصول کی طرح ہم تفرقات  کے مسئلہ اوسط قیمت کی طرف دیکھتے ہیں۔
./calculusConicSectionsParametrizedCurvesPolarCoordinates.tex:4330:ایک قابل تفرق منحنی پر اس نقطہ پر جہاں \عددی{x} کے لحاظ سے بھی \عددی{y} قابل تفرق تفاعل ہو، تفرقات \عددی{\tfrac{\dif x}{\dif t}}، \عددی{\tfrac{\dif y}{\dif t}} اور \عددی{\tfrac{\dif z}{\dif t}} کا تعلق درج ذیل زنجیری قاعدہ دیتا ہے:
./calculusInfiniteSeries.tex:6029:ہم جانتے ہیں کہ اپنے وقفہ ارتکاز کے اندر طاقتی تسلسل کا مجموعہ استمراری تفاعل ہوتا ہے جس کے تفرقات  ہر درجے کے پائے جاتے ہیں۔ لیکن کیا ہم یہی کچھ دوسری رخ بھی کہہ سکتے ہیں؟ یعنی کیا ایسا تفاعل \عددی{f(x)} جس کے وقفہ \عددی{I} ہر رتبہ کے تفرقات پائے جاتے ہوں کو \عددی{ِI} میں طاقتی تسلسل سے ظاہر کرنا ممکن ہو گا؟ اگر ایسا ممکن ہو، تب اس تسلسل کے عددی سر کیا ہوں گے؟
./calculusInfiniteSeries.tex:6110:اگر \عددی{a} پر \عددی{f} کے بلند رتبی تفرقات پائے جاتے ہوں تب ہر پائے جانے والے تفرق کے لئے  اس کی بلند رتبی تخمینی کثیررکنی بھی پائی جائے گی۔  ان کثیررکنیوں کو \عددی{f} کی ٹیلر کثیررکنیاں کہتے ہیں۔
./calculusInfiniteSeries.tex:6195:تفاعل اور اس کے تفرقات درج ذیل ہیں۔
./calculusInfiniteSeries.tex:6220:کثیررکنیاں \عددی{P_{2n}(x)} کوسائن تفاعل پر \عددی{n\to\infty} کرنے سے مرکوز ہوتی ہیں۔ ہم \عددی{x=0} پر کوسائن اور اس کے تفرقات کی قیمتیں جانتے ہوئے کسی بھی فاصلہ پر \عددی{\cos x} کا رویہ جان سکتے ہیں۔
./calculusInfiniteSeries.tex:6242:ہو گا جو ہر \عددی{x} کے لئے مرتکز ہے (جس کا مجموعہ \عددی{0} ہے) لیکن یہ صرف \عددی{x=0} پر \عددی{f(x)} کو مرتکز ہے۔  تفاعل \عددی{y=e^{-1/x^2}} کی استمراری توسیع کی ترسیم (شکل \حوالہ{شکل_مثال_تسلسل_مرکوز_لیکن_تفاعل_کو_نہیں}) صفر پر اتنا (افقی) سیدھا ہے کہ اس کے تمام تفرقات یہاں \عددی{0} کے برابر ہیں۔
./calculusInfiniteSeries.tex:6256:\caption{صفر پر تفاعل کی ترسیم اتنا (افقی) سیدھا ہے  کے یہاں اس کے تمام تفرقات \عددی{0} ہیں۔}
./calculusInfiniteSeries.tex:6482:فرض کریں \عددی{x=a} پر \عددی{f(x)} کے \عددی{n} رتبہ تک تمام تفرقات پائے جاتے ہوں۔ دکھائیں کہ \عددی{x=a} پر \عددی{n} رتبی ٹیلر کثیررکنی اور اس کے ابتدائی \عددی{n} تفرقات کی قیمتیں وہیں ہیں جو \عددی{x=a} پر \عددی{f} اور اس کے ابتدائی \عددی{n} تفرقات کی قیمتیں ہیں۔ 
./calculusInfiniteSeries.tex:6562:اگر \عددی{[a,b]} یا \عددی{[b,a]} پر \عددی{f} اور اس کے \عددی{n} تفرقات \عددی{f'}، \عددی{f''}، \نقطے، \عددی{f^{(n)}} استمراری ہوں اور \عددی{(a,b)} یا \عددی{(b,a)} پر \عددی{f^{(n)}} قابل تفرق ہو تب \عددی{a} اور \عددی{b} کے بیچ ایک ایسا عدد \عددی{c} موجود ہو گا جو درج ذیل کو مطمئن کرتا ہو۔
./calculusInfiniteSeries.tex:6577:اگر وقفہ \عددی{I}، جس میں \عددی{a} پایا جاتا ہو، میں \عددی{f} کے ہر رتبی تفرقات پائے جاتے ہوں، تب ہر مثبت عدد صحیح \عددی{n} اور \عددی{I} میں ہر \عددی{x} کے لئے
./calculusInfiniteSeries.tex:6605:تمام وقفہ \عددی{(-\infty,\infty)} میں اس تفاعل کے ہر رتبی تفرقات پائے جاتے ہیں۔  مساوات \حوالہ{مساوات_تسلسل_ٹیلر_ضمنی_نتیجہ_الف} اور مساوات \حوالہ{مساوات_تسلسل_ٹیلر_ضمنی_نتیجہ_ب} تفاعل \عددی{f(x)=e^x} اور \عددی{a=0} لیتے ہوئے
./calculusInfiniteSeries.tex:6644:سادہ ترین مثال میں ہم \عددی{r=1} لے سکتے ہیں بشرطیکہ \عددی{f} اور اس کے تفرقات کی مقدار کی حد کوئی مستقل \عددی{M} ہو۔ دیگر صورتوں میں ہمیں \عددی{r} کو بھی لینا ہو گا۔ مثال کے طور پر اگر \عددی{f(x)=2\cos (3x)} ہو تب ہر تفرق جزو ضربی \عددی{3} دیگا لہٰذا \عددی{r} کو \عددی{1} سے بڑا منتخب کرنا لازمی ہو گا۔ اس مخصوص مثال میں ہم \عددی{r=3} اور \عددی{M=2} منتخب کر سکتے ہیں۔
./calculusInfiniteSeries.tex:6652:تفاعل اور اس کے تفرقات
./calculusInfiniteSeries.tex:6667:تفاعل \عددی{\sin x} کے تمام تفرقات کی مطلق قیمتیں \عددی{1} یا اس سے کم ہیں لہٰذا ہم اندازہ باقی کے مسئلہ میں \عددی{M=1} اور \عددی{r=1} لیتے ہوئے درج ذیل حاصل کرتے ہیں۔
./calculusInfiniteSeries.tex:6834:اور اس کے ابتدائی \عددی{n} تفرقات، تفاعل \عددی{f} اور اس کے ابتدائی \عددی{n} تفرقات کے موافق ہیں۔ اس میں \عددی{K(x-a)^{n+1}} طرز کا جزو، جہاں \عددی{K} مستقل ہے، شامل کرنے سے یہ موافقت خراب نہیں ہوتی ہے چونکہ  \عددی{x=a} پر ایسا جزو اور اس کے تمام تفرقات صفر ہیں۔نقطہ \عددی{x=a} پر یہ نیا تفاعل
./calculusInfiniteSeries.tex:6838:اور اس کے ابتدائی \عددی{n} تفرقات، \عددی{f} اور اس کے ابتدائی \عددی{n} تفرقات کے ساتھ موافقت رکھیں گے۔
./calculusInfiniteSeries.tex:7402:جس کو \اصطلاح{ثنائی تسلسل}\فرہنگ{ثنائی!تسلسل}\فرہنگ{تسلسل!ثنائی}\حاشیہب{binomial series}\فرہنگ{binomial!series} کہتے ہیں اور جو \عددی{\abs{x}<1} کے لئے مطلق مرتکز ہے۔ یہ تسلسل حاصل کرنے کی خاطر ہم تفاعل اور اس کے تفرقات لکھتے ہیں:
./calculusInfiniteSeries.tex:7651:ہم \عددی{\tan^{-1}x} کا مکلارن تسلسل بلا واسطہ اس لئے دریافت نہیں کرتے ہیں کہ \عددی{\tan^{-1}x} کے بلند رتبی تفرقات کے کلیات پیچیدہ ہوتے ہیں۔ مذکورہ بالا طریقہ زیادہ آسان ہے۔
./calculusIntegration.tex:9:تفاعل کی معلوم قیمتوں میں سے کسی ایک قیمت اور تفاعل کے تفرق \عددی{f(x)} سے تفاعل کا حصول دو قدموں میں ممکن ہے۔ پہلے قدم میں وہ تمام تفاعل حاصل کیے جاتے ہیں جن کا تفرق \عددی{f} ہے۔ ان تفاعل کو \عددی{f} کے الٹ تفرقات کہتے ہیں اور جس کلیہ سے انہیں اخذ کیا جاتا ہے اس کو \عددی{f} کا غیر قطعی  تکمل  کہتے ہیں۔ دوسرے قدم میں تفاعل کی معلوم قیمت استعمال کرتے ہوئے الٹ تفرقات میں سے مخصوص تفاعل منتخب کیا جاتا ہے۔ اس حصہ میں پہلے قدم پر غور کیا جائے گا جبکہ دوسرے قدم پر اگلے حصہ میں غور کیا جائے گا۔
./calculusIntegration.tex:11:اگرچہ تفاعل کے تمام الٹ تفرقات حاصل کرنے والا کلیہ دریافت کرنا  ناممکن نظر آتا ہے، حقیقت میں ایسا نہیں ہے۔ مسئلہ اوسط قیمت (مسئلہ \حوالہ{مسئلہ_استعمال_اوسط_قیمت}) کے  پہلا اور دوسرا ضمنی نتائج کی مدد سے تفاعل کے  ایک الٹ تفرق سے اس کے تمام الٹ تفرقات حاصل کیے جا سکتے ہیں۔ 
./calculusIntegration.tex:19:\عددی{f} کے تمام الٹ تفرقات کا سلسلہ \عددی{x} کے لحاظ سے \عددی{f} کا \اصطلاح{غیر قطعی تکمل}\فرہنگ{تکمل!غیر قطعی}\حاشیہب{indefinite integral}\فرہنگ{integral!indefinite} ہو گا جس کو درج ذیل سے ظاہر کیا جاتا ہے۔
./calculusIntegration.tex:39:\عددی{2x} کا الٹ تفرق \عددی{x^2} ہے اور \عددی{C} تکمل کا مستقل ہے۔کلیہ \عددی{x^2+C} تفاعل \عددی{2x} کے تمام تفرقات دیتا ہے۔یوں \عددی{x^2+1}، \عددی{x^2-\pi} اور \عددی{x^2+\sqrt{2}} تفاعل \عددی{2x} کے ممکنہ الٹ تفرق ہیں۔ آپ ان کا تفرق لے کر تصدیق کر سکتے ہیں۔
./calculusIntegration.tex:43:ہم عموماً تفرق کے کلیات سے الٹ تفرقات کے کلیات اخذ کرتے ہیں۔جدول \حوالہ{جدول_تکمل_کلیات_الف} میں غیر قطعی تکملات کے سامنے موزوں تفرقی کلیات کو الٹ لکھا گیا ہے۔ 
./calculusIntegration.tex:102:\جزوحصہء{الٹ تفرقات کے قواعد}
./calculusIntegration.tex:103:ہم الٹ تفرقات کے بارے میں درج ذیل جانتے ہیں۔
./calculusMultipleIntegrals.tex:6696:اس کا جواب: اگر \عددی{g}، \عددی{h} اور \عددی{f} کے جزوی تفرقات استمراری ہوں اور \عددی{J(u,v)} (جس پر جلد تبصرہ کیا جائے گا) صرف  تنہا نقطوں پر صفر ہو (اگر صفر ہو بھی)  تب درج ذیل ہو گا۔ 
./calculusMultipleIntegrals.tex:6720:سے بھی ظاہر کیا جاتا ہے جو ہمیں یاد دلاتا ہے کہ \عددی{x} اور \عددی{y} کی جزوی تفرقات سے یعقوبی (مساوات \حوالہ{مساوات_بالکثرت_یعقوبی_تعریف})  حاصل ہوتا ہے۔مساوات \حوالہ{مساوات_بالکثرت_یعقوبی_بدل} کی استخراج آپ کو اعلیٰ احصاء کے نصاب میں ملے گی جس کو یہاں پیش نہیں کیا جائے گا۔
./calculusMultipleIntegrals.tex:7012:تصور کیا جا سکتا ہے۔اگر \عددی{g}، \عددی{h}  اور \عددی{k} کے اول جزوی تفرقات استمراری ہوں تب \عددی{D} پر \عددی{F} کے تکمل کا \عددی{G} پر \عددی{H(u,v,w)} کے تکمل کے ساتھ تعلق درج ذیل مساوات دیگی۔
./calculusMultivariableFunctionsAndPartialDerivatives.tex:1:\باب{کثیر المتغیر تفاعل اور جزوی تفرقات}\شناخت{باب_کثیر_المتغیر_تفاعل_اور_جزوی_تفرقات}
./calculusMultivariableFunctionsAndPartialDerivatives.tex:3:سائنس میں دو یا دو سے زائد غیر تابع  متغیرات کے تفاعل    ایک متغیر کے تفاعل سے زیادہ کثرت سے پائے جاتے ہیں اور ان کی علم   احصاء  زیادہ  عمدہ   ہوتی  ہے۔زیادہ متغیرات ایک دوسرے پر زیادہ طریقوں سے اثر انداز ہو سکتے ہیں جس کی وجہ سے ان کے تفرقات    مختلف اور زیادہ دلچسپ  صورتیں اختیار  کر سکتے ہیں۔ ان کے تکملات  زیادہ اقسام کے عملی  مسائل میں کام آتے ہیں۔ احتمال،  سیالی حرکیات، اور برقیات، وغیرہ،    پر غور کے دوران  ایک سے زائد متغیرات  کے تفاعل قدرتی طور پر رونما ہوتے ہیں۔ان تفاعل کی  ریاضیات، سائنس کی عظیم  کامیابیوں میں  سے ایک ہے۔
./calculusMultivariableFunctionsAndPartialDerivatives.tex:1960:\حصہ{جزوی تفرقات}\شناخت{حصہ_کثیر_متغیر_جزوی_تفرق}
./calculusMultivariableFunctionsAndPartialDerivatives.tex:1961:جب ہم   ماسوائے ایک غیر تابع متغیر  کے باقی تمام  کو برقرار رکھیں اور اس ایک متغیر کے لحاظ سے تفاعل کا تفرق لیں تو ہمیں \قول{جزوی} تفرق حاصل ہوتا ہے۔ اس حصہ میں دکھایا جائے گا کہ جزوی تفرقات کیسے  پائے جاتے ہیں اور واحد متغیر کے تفاعل کے تفرق کے قواعد بروئے کار لاتے ہوئے   جزوی تفرقات   کی قیمت کے حصول کے بارے میں بتایا جائے گا۔
./calculusMultivariableFunctionsAndPartialDerivatives.tex:2010:جہاں جزوی تفرق  کو ایک تفاعل تصور کرنا مقصود  ہو وہاں درج ذیل علامت استعمال کیے جائیں گے، جہاں  \عددی{x} لے لحاظ سے \عددی{f} (یا \عددی{z}) کے جزوی تفرقات لیے گئے ہیں۔
./calculusMultivariableFunctionsAndPartialDerivatives.tex:2116:\ترچھا{جزوی تفرق}\quad کمپیوٹر آپ کو حساب میں کئی بعد  تک مدد فراہم کر سکتا ہے۔ آپ ایک غیر تابع متغیر کے علاوہ تمام متغیرات کی قیمتیں فراہم کر کے واحد  متغیر کے لحاظ سے جزوی تفرق معلوم کر  کے  ترسیم کر سکتے ہیں۔ جزوی تفرق اور سادہ تفرق کے لئے کمپیوٹر کی زبان میں عموماً  ایک  سی  اصطلاح استعمال کی جاتی ہے۔  جزوی تفرقات کے حصول میں کمپیوٹر ضرور استعمال کریں۔
./calculusMultivariableFunctionsAndPartialDerivatives.tex:2256:نقطہ \عددی{(0,0)} پر غیر استمراری ہے (شکل \حوالہ{شکل_مثال_کثیرالمتغیر_ٹکڑوں_میں_استمراری})۔ لکیر \عددی{y=x} پر چلتے ہوئے  نقطہ  \عددی{(x,y)} کا  \عددی{(x_0,y_0)} تک پہنچنے سے \عددی{f} کی حد \عددی{0} حاصل ہوتی  ہے جبکہ \عددی{f(0,0)=1} ہے۔   نقطہ \عددی{(0,0)} پر \عددی{f} کے جزوی تفرقات \عددی{f_x}، \عددی{f_y}، جو شکل میں خط \عددی{L_1} اور \عددی{L_2} کی ڈھلوان ہیں،  دونوں موجود ہیں۔
./calculusMultivariableFunctionsAndPartialDerivatives.tex:2260:\جزوحصہء{دو رتبی جزوی  تفرقات}
./calculusMultivariableFunctionsAndPartialDerivatives.tex:2261:تفاعل \عددی{f(x,y)} کو دو بار تفرق کرنے سے ہمیں  اس تفاعل کا دو رتبی تفرق ملتا ہے۔ان تفرقات کو عموماً درج ذیل سے ظاہر کیا جاتا ہے۔
./calculusMultivariableFunctionsAndPartialDerivatives.tex:2295:آپ نے مثال \حوالہ{مثال_کثیرالمتغیر_مدغم_دو_رتبی} میں دھیان دیا ہو گا کہ    مدغم      دو رتبی جزوی تفرقات
./calculusMultivariableFunctionsAndPartialDerivatives.tex:2302:اگر   ایک  کھلے  خطہ میں،  جس میں نقطہ \عددی{(a,b)} پایا جاتا ہو، \عددی{f(x,y)} اور اس کے جزوی تفرقات \عددی{f_x}، \عددی{f_y}، \عددی{f_{xy}} اور \عددی{f_{yx}} معین ہوں اور    \عددی{(a,b)}  پر یہ تمام  استمراری ہوں تب درج ذیل ہو گا۔
./calculusMultivariableFunctionsAndPartialDerivatives.tex:2328:\جزوحصہء{مزید بلند رتبہ کے جزوی تفرقات}
./calculusMultivariableFunctionsAndPartialDerivatives.tex:2329:عملی استعمال میں یک رتبی اور دو رتبی جزوی تفرقات   زیادہ کثرت سے پائے جاتے ہیں لہٰذا ہمیں عموماً انہیں سے واسطہ ہو گا ۔جہاں تک تفاعل کے بلند تفرقات  کی بات ہے، ہم ایک تفاعل کا تفرق  جتنی بار چاہیں لیں سکتے ہیں بشرطیکہ ایسے تفرقات موجود ہوں۔ یوں ہم تین رتبی اور چار رتبی تفرقات لے سکتے ہیں جنہیں درج ذیل علامتوں  کی طرز    پر  ظاہر کیا جائے گا۔
./calculusMultivariableFunctionsAndPartialDerivatives.tex:2334:دو رتبی تفرق کی طرح،   تفرق کی ترتیب غیر اہم ہے جب تک تمام  تفرقات استمراری ہوں۔
./calculusMultivariableFunctionsAndPartialDerivatives.tex:2645:سوال \حوالہ{سوال_کثیرالمتغیر_تمام_جزوی_تلاش_الف} تا سوال \حوالہ{سوال_کثیرالمتغیر_تمام_جزوی_تلاش_ب} میں تفاعل کے تمام دو رتبی جزوی تفرقات تلاش کریں۔
./calculusMultivariableFunctionsAndPartialDerivatives.tex:2695:\موٹا{مدغم جزوی تفرقات}\\
./calculusMultivariableFunctionsAndPartialDerivatives.tex:2800:\موٹا{خفی جزوی تفرقات}\\
./calculusMultivariableFunctionsAndPartialDerivatives.tex:2981:\حصہ{تفرق  پذیری، خط بندی، اور تفرقات}\شناخت{حصہ_کثیرالمتغیر_تفرق_پذیری_خط_بندی_تفرقات}
./calculusMultivariableFunctionsAndPartialDerivatives.tex:2992:فرض کریں    پورا کھلا  خطہ \عددی{R} میں، جس میں نقطہ \عددی{(x_0,y_0)} پایا جاتا ہو،     \عددی{f(x,y)} کے    جزوی  اول تفرقات  معین  ہیں اور \عددی{(x_0,y_0)} پر \عددی{f_x} اور \عددی{f_y} استمراری ہیں۔تب نقطہ \عددی{(x_0,y_0)} کو \عددی{R} میں    دوسری جگہ  \عددی{(x_0+\Delta x,y_0+\Delta y)} منتقل کرنے سے \عددی{f} میں رونما ہونے والی تبدیلی
./calculusMultivariableFunctionsAndPartialDerivatives.tex:3010:اس تعریف کی روشنی میں ہمیں مسئلہ \حوالہ{مسئلہ_کثیرالمتغیر_بڑھوتری} کا  ضمنی نتیجہ ملتا ہے جس کے تحت  جس تفاعل کے   جزوی اول   تفرقات استمراری ہوں وہ تفاعل قابل تفرق ہو گا۔
./calculusMultivariableFunctionsAndPartialDerivatives.tex:3013:اگر پورے کھلا وقفہ \عددی{R} میں تفاعل \عددی{f(x,y)} کے جزوی تفرقات \عددی{f_x} اور \عددی{f_x} استمراری ہوں تب \عددی{R} کے ہر نقطہ پر \عددی{f} تفرق پذیر ہو گا۔
./calculusMultivariableFunctionsAndPartialDerivatives.tex:3113:تخمین \عددی{f(x,y)\approx L(x,y)}  میں خلل کی تلاش میں ہم \عددی{f} کے دو  رتبی  جزوی تفرقات استعمال کرتے ہیں۔ فرض کریں ایک  کھلا  سلسلہ  میں \عددی{f} کے یک رتبی اور دو رتبی جزوی تفرقات استمراری ہوں اور اس سلسلہ   میں ایک  مستطیل خطہ \عددی{R}   جس کا مرکز \عددی{(x_0,y_0)} ہو پایا جاتا ہو۔ اس مستطیل خطہ کو  درج ذیل عدم مساوات  ظاہر کرتے ہیں (شکل \حوالہ{شکل_کثیرالمتغیر_تخمین_درستگی})۔
./calculusMultivariableFunctionsAndPartialDerivatives.tex:3117:چونکہ \عددی{R} بند اور محدود ہے لہٰذا \عددی{R} میں    تمام دو رتبی جزوی تفرقات  کی مطلق   زیادہ سے زیادہ قیمتیں ہوں گی۔ اگر ان میں \عددی{B} سب سے بڑی قیمت ہو تب، جیسا   آگے حصہ \حوالہ{حصہ_کثیر_کلیہ_ٹیلر} میں سمجھایا گیا ہے،پورے \عددی{R} میں   معیاری خطی تخمین میں خلل \عددی{E(x,y)=f(x,y)-L(x,y)}  درج ذیل عدم مساوات کو مطمئن کرے گا۔ 
./calculusMultivariableFunctionsAndPartialDerivatives.tex:3129:اگر   ایک  کھلا  سلسلہ  میں \عددی{f} کے یک رتبی اور دو رتبی جزوی تفرقات استمراری ہوں اور اس سلسلہ   میں ایک  مستطیل خطہ \عددی{R}   جس کا مرکز \عددی{(x_0,y_0)} ہو پایا جاتا ہو  اور \عددی{R} پر \عددی{\abs{f_{xx}}}، \عددی{\abs{f_{yy}}} اور \عددی{\abs{f_{xy}}}  کی بالائی حد  بندی \عددی{M} ہو   تب  \عددی{R}  پر \عددی{f(x,y)} کا متبادل 
./calculusMultivariableFunctionsAndPartialDerivatives.tex:3172:فرض  کریں ہم   نقطہ \عددی{(x_0,y_0)}   پر قابل تفرق تفاعل  \عددی{f(x,y)} اور اس کے یک رتبی تفرقات  کی قیمتیں جانتے ہیں اور ہم قریبی نقطہ \عددی{(x_0+\Delta x,y_0+\Delta y)}  پر   منتقل ہونے سے \عددی{f} کی قیمت میں تبدیلی  جاننا چاہتے  ہیں۔ اگر \عددی{\Delta x} اور \عددی{\Delta y} چھوٹے ہوں تب\عددی{(x_0,y_0)} پر  \عددی{f} اور اس کی خط بندی کی  قیمت  میں تبدیلی تقریباً ایک دوسرے جیسی ہو گی لہٰذا \عددی{L} کی تبدیلی سے ہمیں عملاً  \عددی{f} کی تبدیلی حاصل ہو گی۔
./calculusMultivariableFunctionsAndPartialDerivatives.tex:3375:فرض کریں بند  ٹھوس مستطیل  \عددی{R} کا مرکز \عددی{N_0} ہے۔ یہ مستطیل ایسے خطہ میں پایا جاتا ہے جہاں \عددی{f} کے دو رتبی جزوی تفرقات استمراری ہیں۔ مزید فرض کریں کہ پورے \عددی{R} پر \عددی{\abs{f_{xx}}}،  \عددی{\abs{f_{yy}}}،  \عددی{\abs{f_{zz}}}،  \عددی{\abs{f_{xy}}}،  \عددی{\abs{f_{xz}}} اور  \عددی{\abs{f_{yz}}} کی قیمتیں \عددی{M} کے برابر یا اس سے کم ہیں۔تب پورے \عددی{R} میں   \عددی{f} کی تخمین \عددی{L} میں \اصطلاح{خلل}\فرہنگ{خلل}\حاشیہب{error}\فرہنگ{error} \عددی{ٰE(x,y,z)=f(x,y,z)-L(x,y,z)}کی بالائی حد بندی   درج ذیل عدم مساوات دے گی۔
./calculusMultivariableFunctionsAndPartialDerivatives.tex:3380:اگر \عددی{f} کے  دو رتبی جزوی تفرقات استمراری ہوں اور \عددی{x}، \عددی{y}، \عددی{z} کی قیمتیں  چھوٹی تبدیلیوں \عددی{\dif x}، \عددی{\dif y}، \عددی{\dif z} کی وجہ سے  \عددی{x_0}، \عددی{y_0}، \عددی{z_0} سے تبدیل ہو جائیں ، تب کل  تفریق
./calculusMultivariableFunctionsAndPartialDerivatives.tex:3405:اسی طرح پہلے دو رتبی جزوی تفرقات حاصل کرتے ہیں۔
./calculusMultivariableFunctionsAndPartialDerivatives.tex:3863:کیا پورے  کھلا خطہ \عددی{R} میں استمراری دو رتبی جزوی تفرقات والا   تفاعل \عددی{f(x,y)} خطہ \عددی{R} میں لازماً استمراری ہو گا؟ اپنے جواب کی وجہ پیش کریں۔
./calculusMultivariableFunctionsAndPartialDerivatives.tex:3873:کیا پورے کھلا خطہ \عددی{R} میں استمراری دو رتبی جزوی تفرقات والے ا   تفاعل \عددی{f(x,y)}کے  خطہ \عددی{R} میں لازماً استمراری یک رتبی جزوی تفرقات پائے جائیں گے؟ اپنے جواب کی وجہ پیش کریں۔
./calculusMultivariableFunctionsAndPartialDerivatives.tex:3922:تفرق \عددی{\tfrac{\dif w}{\dif t}} میں حقیقی غیر تابع متغیر \عددی{t} اور تابع متغیر \عددی{w} ہے جبکہ \عددی{x} اور \عددی{y}  درمیانی متغیرات ہیں جنہیں \عددی{t} قابو کرتا ہے۔زنجیری قاعدہ کا  درج ذیل روپ ہمیں  مساوات \حوالہ{مساوات_کثیرالمتغیر_زنجیری_دو_متغیرات_الف} میں مختلف تفرقات کے حصول کا صحیح طریقہ دیتا ہے۔
./calculusMultivariableFunctionsAndPartialDerivatives.tex:3926:یوں \عددی{\tfrac{\dif x}{\dif t}} اور \عددی{\tfrac{\dif y}{\dif t}} نقطہ \عددی{t_0} پر حاصل کیے جائیں گے۔ حقیقی غیر تابع متغیر کی قیمت  \عددی{t_0}،  درمیانی متغیرات \عددی{x} اور \عددی{y} کی   \عددی{x_0} اور \عددی{y_0} قیمتیں  تعین کرتا ہے۔ جزوی تفرقات \عددی{\tfrac{\partial w}{\partial x}} اور \عددی{\tfrac{\partial w}{\partial y}} نقطہ \عددی{(x_0,y_0)} پر حاصل کیے جاتے ہیں۔ 
./calculusMultivariableFunctionsAndPartialDerivatives.tex:3928:زنجیری  قاعدہ کو شکل \حوالہ{شکل_کثیر_المتغیر_زنجیری_قاعدہ_یاد_رکھنا} کی  مدد سے یاد رکھنا زیادہ آسان  ہو گا۔اس طرح شکل  کو  \اصطلاح{شکل شجرہ}\فرہنگ{شکل شجرہ}\حاشیہب{tree diagram}\فرہنگ{tree diagram}  کہتے ہیں۔  آپ شکل شجرہ سے  دیکھ سکتے ہیں کہ جب \عددی{t=t_0} ہو تفرقات \عددی{\tfrac{\dif x}{\dif t}} اور \عددی{\tfrac{\dif y}{\dif t}} کی قیمتیں   \عددی{t_0} پر حاصل کیا جاتا ہے۔ اب  قابل تفرق تفاعل \عددی{x} کے لئے \عددی{x_0} کی قیمت \عددی{t_0} تعین کرتا ہے اور قابل تفرق تفاعل \عددی{y} کے لئے \عددی{y_0}  کی قیمت \عددی{t_0} تعین کرتا ہے۔ جزوی تفرقات \عددی{\tfrac{\partial w}{\partial x}}، \عددی{\tfrac{\partial w}{\partial y}} (جو از خود \عددی{x} اور \عددی{y} کے تفاعل ہیں)  کی قیمتیں نقطہ \عددی{N_0(x_0,y_0)} پر حاصل کی جاتی ہیں جو \عددی{t_0} کا مطابقتی نقطہ ہے۔ حقیقی  غیر تابع متغیر \عددی{t} ہے جبکہ \عددی{x} اور \عددی{y} درمیانے متغیرات اور \عددی{w}  تابع متغیر ہے۔
./calculusMultivariableFunctionsAndPartialDerivatives.tex:3949:زنجیری قاعدہ یاد رکھنے کی خاطر اس شکل کو ذہن میں رکھیں۔اب \عددی{\tfrac{\dif w}{\dif t}}  معلوم کرنے کی خاطر \عددی{w} سے شروع ہو کر \عددی{t} تک باری باری دونوں راہ پر  چل  کر  تفرقات کا  حاصل ضرب لیں ۔  آخر میں دونوں راہ کے حاصل ضرب کا  مجموعہ لیں۔ 
./calculusMultivariableFunctionsAndPartialDerivatives.tex:3969:آپ نے دیکھا کہ ہم نے \عددی{x=\cos t} اور \عددی{y=\sin t} کو جزوی تفرقات \عددی{\tfrac{\partial w}{\partial x}} اور \عددی{\tfrac{\partial w}{\partial y}} میں پر کیا (شکل \حوالہ{شکل_مثال_کثیرالمتغیر_زنجیری_حاصل_ضرب_تفاعل})۔ یوں حاصل تفرق \عددی{\tfrac{\dif w}{\dif t}}  کا اظہار غیر تابع متغیر \عددی{t} کی صورت میں   کیا جاتا ہے (جس میں  درمیا نے    متغیرات \عددی{x} اور \عددی{y} نہیں پائے  جاتے ہیں۔ )
./calculusMultivariableFunctionsAndPartialDerivatives.tex:4009:\caption{یہاں \عددی{w} سے \عددی{t} تک تین راستے ہیں۔اب بھی \عددی{\tfrac{\dif w}{\dif t}} حاصل کرنے کی خاطر ہر راہ پر چلتے ہوئے تفرقات کا ضرب لے کر تمام    کا مجموعہ لیں۔}
./calculusMultivariableFunctionsAndPartialDerivatives.tex:4070: تصور کر سکتے ہیں۔ موزوں حالات میں  \عددی{r} اور \عددی{s} دونوں کے لحاظ سے \عددی{w} کے جزوی تفرقات موجود ہوں گے جنہیں درج ذیل طریقہ سے حاصل کیا جا سکتا ہے۔
./calculusMultivariableFunctionsAndPartialDerivatives.tex:4073:فرض کریں \عددی{w=f(x,y,z)}، \عددی{x=g(r,s)}، \عددی{y=h(r,s)} اور \عددی{z=k(r,s)} ہوں۔ اگر چاروں تفاعل قابل تفرق ہوں، تب \عددی{r} اور \عددی{s} کے لحاظ سے \عددی{w} کے جزوی تفرقات قابل  پائے جائیں گے جنہیں درج ذیل مساوات سے حاصل کیا جا سکتا ہے۔
./calculusMultivariableFunctionsAndPartialDerivatives.tex:4310:فرض کریں \عددی{f(x,y,\cdots,v)} غیر تابع (متناہی تعداد کے)  متغیرات \عددی{x,y,\cdots,v}  کا قابل تفرق تفاعل ہے اور \عددی{x,y,\cdots,v}  (ایک دوسرے متناہی تعداد کے)  متغیرات \عددی{p,q,\cdots,t} کے قابل تفرق تفاعل ہیں۔ تب \عددی{w} متغیرات \عددی{p,q,\cdots,t} کا قابل تفرق تفاعل ہو گا اور ان متغیرات کے لحاظ سے \عددی{w} کے جزوی تفرقات  کی  صورت درج ذیل ہو گی۔
./calculusMultivariableFunctionsAndPartialDerivatives.tex:4318:\underbrace{\big(\frac{\partial w}{\partial x},\frac{\partial w}{\partial y},\cdot,\frac{\partial w}{\partial v}\big)}_{\text{\RL{درمیانی  متغیرات کے لحاظ سے \عددی{w} کے تفرقات}}}\quad \text{اور}\quad \underbrace{\big(\frac{\partial x}{\partial p},\frac{\partial y}{\partial p},\cdot,\frac{\partial v}{\partial p}\big)}_{\text{\RL{منتخب غیر تابع متغیر کے لحاظ سے درمیانی متغیرات کے تفرقات}}}
./calculusMultivariableFunctionsAndPartialDerivatives.tex:4596:\موٹا{منتخب جزوی تفرقات}\\
./calculusMultivariableFunctionsAndPartialDerivatives.tex:4693:فرض کریں پیچدار منحنی \عددی{x=\cos t,\,y=\sin t,\, z=t} پر  پائے جانے والے نقطوں پر تفاعل \عددی{f(x,y,z)}  کے جزوی تفرقات درج ذیل ہیں۔
./calculusMultivariableFunctionsAndPartialDerivatives.tex:4717:تفرقات \عددی{\tfrac{\dif T}{\dif t}} اور \عددی{\tfrac{\dif^{\,2}T}{\dif t^2}} کو دیکھ کر بتائیں کہ اس دائرہ پر کہاں زیادہ سے زیادہ اور کہاں کم سے کم درجہ حرارت ہو گا۔
./calculusMultivariableFunctionsAndPartialDerivatives.tex:4744:تفرقات \عددی{\tfrac{\dif T}{\dif t}} اور \عددی{\tfrac{\dif^{\,2}T}{\dif t^2}} کو دیکھ کر بتائیں  ترخیم پر  کہاں زیادہ سے زیادہ  اور  کہاں کم سے کم    \عددی{T} ہو گا۔
./calculusMultivariableFunctionsAndPartialDerivatives.tex:4750:\موٹا{تکملات کے تفرقات}\\
./calculusMultivariableFunctionsAndPartialDerivatives.tex:4784:اب تک تفاعل، مثلاً  \عددی{w=f(x,y)}،    کے جزوی تفرقات تلاش کرتے ہوئے  ہم \عددی{x} اور \عددی{y} کو بالکل آزاد  غیر تابع متغیرات تصور کرتے رہے ہیں، اگرچہ  عملی زندگی میں ضروری نہیں کہ ایسا ہو۔ مثال کے طور پر ہم گیس کی اندرونی  توانائی  \عددی{U} کو دباو \عددی{P}، حجم \عددی{H} اور حرارت \عددی{T} کا تفاعل \عددی{U=f(P,H,T)} لکھ سکتے ہیں۔ اگر گیس کے انفرادی مالیکیول ایک دوسرے  پر اثر انداز نہ ہوں تب  \عددی{P}، \عددی{H} اور \عددی{T}  مثالی گیس کے    قانون
./calculusMultivariableFunctionsAndPartialDerivatives.tex:4788: کو مطمئن کریں گے لہٰذا یہ متغیرات بالکل آزاد ہرگز نہیں ہوں گے۔ایسی صورت میں جزوی تفرقات کی تلاش  پیچیدہ ثابت ہوتے ہیں۔بہر حال ان سے نمٹنا ضروری ہے۔ 
./calculusMultivariableFunctionsAndPartialDerivatives.tex:4791:اگر تفاعل \عددی{w=f(x,y,z)} کے متغیرات کسی تعلق، مثلاً \عددی{z=x^2+y^2}، کے پابند ہوں تب \عددی{f} کے جزوی تفرقات کی جیومیٹریائی  معنی اور عددی قیمت  اس پر منحصر ہوں گے کہ کن متغیرات کو غیر تابع اور کن کو تابع متغیرات تصور کیا جاتا ہے۔ اس انتخاب کے اثرات کو دیکھنے کی خاطر آئیں \عددی{w=x^2+y^2+z^2}  اور \عددی{z=x^2+y^2} کی صورت میں \عددی{\tfrac{\partial w}{\partial x}} تلاش کریں۔ 
./calculusMultivariableFunctionsAndPartialDerivatives.tex:4859:\caption{نقطہ \عددی{N} کو پابند کرنے سے جزوی تفرقات کے مختلف نتائج حاصل ہوں گے۔}
./calculusMultivariableFunctionsAndPartialDerivatives.tex:4943:اس کلیہ کے  دائیں ہاتھ میں \عددی{w} کے جزوی تفرقات کو  \عددی{w=x^2+y-z+\sin t} کے لئے حاصل کر کے  اس کلیہ میں پر کرتے ہوئے  درج ذیل حاصل کرتے ہیں۔
./calculusMultivariableFunctionsAndPartialDerivatives.tex:4950:باقی جزوی تفرقات حاصل کرنے کے لئے  ہم متغیرات کی غیر تابعیت اور تابعیت  بروئے کار لاتے ہیں۔ ہم مساوات \حوالہ{مساوات_کثیرالمتغیر_تیر_دار} سے دیکھتے ہیں کہ \عددی{x}، \عددی{y} اور \عددی{z} غیر تابع  ہیں اور \عددی{t=x+y}  ہے۔ یوں
./calculusMultivariableFunctionsAndPartialDerivatives.tex:4963:\موٹا{پابند متغیرات کے تفاعل کے جزوی تفرقات}\\
./calculusMultivariableFunctionsAndPartialDerivatives.tex:4964:سوال \حوالہ{سوال_کثیرالمتغیر_پابند_تفرق_الف} تا سوال \حوالہ{سوال_کثیرالمتغیر_پابند_تفرق_ب} میں  تیر دار اشکال سے شروع کرتے ہوئے دیے تفرقات تلاش کریں۔
./calculusMultivariableFunctionsAndPartialDerivatives.tex:5018:فرض کریں \عددی{x^2+y^2=r^2} اور \عددی{x=r\cos\theta} ہیں۔ جزوی تفرقات  \عددی{\big(\tfrac{\partial x}{\partial r}\big)_{\theta}} اور 
./calculusMultivariableFunctionsAndPartialDerivatives.tex:5036:\موٹا{ بغیر کسی مخصوص کلیہ جزوی تفرقات کا حصول}\\
./calculusMultivariableFunctionsAndPartialDerivatives.tex:5042:ہو گا۔اس حقیقت  کی تصدیق کریں۔ (اشارہ:تمام تفرقات کو باضابطہ جزوی تفرقات \عددی{\tfrac{\partial f}{\partial x}}،  \عددی{\tfrac{\partial f}{\partial y}} اور \عددی{\tfrac{\partial f}{\partial z}} کی صورت میں لکھیں۔)
./calculusMultivariableFunctionsAndPartialDerivatives.tex:5075:\حصہ{ رخی تفرقات، سمتیہ ڈھلوان، اور مماسی سطحیں}\شناخت{حصہ_کثیرالمتغیر_رخی_تفرقات_سمتیہ_ڈھلوان_مماسی_سطحیں}
./calculusMultivariableFunctionsAndPartialDerivatives.tex:5080:نقطہ \عددی{N_0(x_0,y_0)=N_0(g(t_0),h(t_0))}  پر یہ مساوات بڑھتی \عددی{t} کے لحاظ سے \عددی{f} کی شرح تبدیلی دیتی ہے،  جو دیگر  چیزوں کے ساتھ  منحنی پر چلنے کے رخ پر بھی منحصر ہے۔ یہ  مشاہدہ   اس صورت خصوصاً     اہم  ہو  گا  جب یہ    منحنی ایک سیدھی  لکیر ہو  اور نقطہ  \عددی{N_0} سے منحنی پر اکائی سمتیہ \عددی{\uu} کے رخ  چلتے ہوئے  مقدار معلوم لمبائی قوس   \عددی{t}  ہو۔چونکہ تب  \عددی{\uu} کے رخ \عددی{f} کے دائرہ کار میں  فاصلہ کے لحاظ  سے   \عددی{f} کی شرح تبدیلی  \عددی{\tfrac{\dif f}{\dif t}} ہو گی۔ ہم    \عددی{\uu}  تبدیل کرتے ہوئے نقطہ \عددی{N_0} پر،     فاصلہ کے لحاظ سے \عددی{f} کی مختلف رخ میں  شرح تبدیلی    دریافت کر سکتے ہیں۔   ان  \اصطلاح{رخی تفرقات}\فرہنگ{تفرق!رخی}\حاشیہب{directional derivative}\فرہنگ{derivative!directional} کی  سائنس،  انجینئری اور ریاضیات   میں کارآمد  تشریحات  کی جاتی ہیں۔ اس حصہ میں ان کی قیمت دریافت کرنے کا کلیہ اخذ کیا جائے گا جس کے بعد فضا میں سطحوں کی مماسی سطحیں اور عمودی سطحیں تلاش کی جائیں گی۔
./calculusMultivariableFunctionsAndPartialDerivatives.tex:5082:\جزوحصہء{مستوی میں رخی تفرقات}
./calculusMultivariableFunctionsAndPartialDerivatives.tex:5136:دھیان رہے کہ جب \عددی{\uu=\ai} ہو، \عددی{N_0} پر  رخی تفرق \عددی{\tfrac{\partial f}{\partial x}} ہو گا جس کی قیمت \عددی{(x_0,y_0)} پر حاصل کی جائے گی۔ اسی طرح  جب   \عددی{\uu=\aj} ہو، \عددی{N_0} پر  رخی تفرق \عددی{\tfrac{\partial f}{\partial y}} ہو گا جس کی قیمت \عددی{(x_0,y_0)} پر حاصل کی جائے گی۔ رخی تفرق ان دو جزوی تفرقات کو عمومی بناتا ہے۔ ہم اب \عددی{\ai} اور \عددی{\aj} کے علاوہ کسی بھی رخ  \عددی{\uu}،    تفاعل \عددی{f} کی تبدیلی شرح  جان سکتے ہیں۔
./calculusMultivariableFunctionsAndPartialDerivatives.tex:5164:اگر \عددی{N_0(x_0,y_0)} پر \عددی{f(x,y)} کے جزوی تفرقات معین ہوں تب \عددی{N_0} پر \عددی{\uu} کے رخ \عددی{f} کا تفرق،   \عددی{N_0} پر \عددی{f} کی ڈھلوان اور \عددی{\uu} کا  غیر سمتی  ضرب    ہو گا:
./calculusMultivariableFunctionsAndPartialDerivatives.tex:5179:نقطہ \عددی{(2,0)} پر \عددی{f} کے جزوی تفرقات
./calculusMultivariableFunctionsAndPartialDerivatives.tex:5208:\جزوحصہء{رخی تفرقات  کے خواص}
./calculusMultivariableFunctionsAndPartialDerivatives.tex:5386:نقطہ \عددی{N_0} پر جزوی تفرقات
./calculusMultivariableFunctionsAndPartialDerivatives.tex:5560:نقطہ \عددی{N_0(x_0,y_0,z_0)} پر     سطح \عددی{z=f(x,y)}  کے مماسی مستوی کی مساوات تلاش کریں۔ اس نقطہ پر  \عددی{z_0=f(x_0,y_0)}  ہو گا۔ ہم دیکھتے ہیں کہ مساوات \عددی{z=f(x,y)}  کو \عددی{f(x,y)-z=0} لکھا  جا سکتا ہے۔ یوں سطح \عددی{z=f(x,y)}  درحقیقت تفاعل \عددی{F(x,y,z)=f(x,y)-z} کا صفر ہم قد مستوی ہو گا۔تفاعل \عددی{F}  کے جزوی تفرقات
./calculusMultivariableFunctionsAndPartialDerivatives.tex:5579:ہم  جزوی تفرقات معلوم کر کے مساوات \حوالہ{مساوات_کثیرالمتغیر_مماسی_مستوی} استعمال کریں گے:
./calculusMultivariableFunctionsAndPartialDerivatives.tex:6109:نقطہ \عددی{N(1,1)} پر کن دو رخ تفاعل \عددی{f(x,y)=\tfrac{x^2-y^2}{x^2+y^2}} کے تفرقات  صفر ہوں گے؟
./calculusMultivariableFunctionsAndPartialDerivatives.tex:6213:نقاط  \عددی{t=-\pi/4,0,\pi/4} پر پیچ دار  منحنی \عددی{\kvec{r}(t)=(\cos t)\ai+(\sin t)\aj+t\ak}  کے اکائی مماسی سمتیات کے رخ     تفاعل \عددی{f(x,y,z)=x^2+y^2+z^2} کے تفرقات   تلاش کریں۔ یہ تفاعل مبدا سے پیچدار منحنی کے نقطہ \عددی{N(x,y,z)} کے  فاصلہ کا مربع  دیتا ہے۔اس منحنی پر چلتے ہوئے   تفرق   \عددی{t} کے لحاظ سے فاصلے کے مربع کی شرح تبدیلی دے گا۔
./calculusMultivariableFunctionsAndPartialDerivatives.tex:6250:\ابتدا{سوال}\ترچھا{رخی تفرقات اور غیر سمتی اجزاء}\\
./calculusMultivariableFunctionsAndPartialDerivatives.tex:6251:نقطہ \عددی{N_0} پر اکائی سمتیہ \عددی{\uu}  کے رخ قابل تفرق تفاعل \عددی{f(x,y,z)}  کے تفرقات  کا  \عددی{\uu} کے رخ \عددی{(\nabla f)_{N_0}}  کے غیر سمتی اجزاء کے ساتھ کیا  تعلق ہے؟ اپنے جواب کی وجہ پیش کریں۔
./calculusMultivariableFunctionsAndPartialDerivatives.tex:6254:\ابتدا{سوال}\ترچھا{رخی تفرقات اور جزوی تفرقات}\\
./calculusMultivariableFunctionsAndPartialDerivatives.tex:6255:فرض کریں کہ \عددی{f(x,y,z)} کے مطلوبہ تفرقات موجود ہیں۔ تفرقات \عددی{D_{\ai}f}، \عددی{D_{\aj}f} اور \عددی{D_{\ak}f} کا تفرقات \عددی{f_x}، \عددی{f_y} اور \عددی{f_z} کے ساتھ کیا تعلق ہے؟ اپنے جواب کی وجہ پیش کریں۔
./calculusMultivariableFunctionsAndPartialDerivatives.tex:6284:مستوی \عددی{xy} میں محدود   بند خطہ میں استمراری تفاعل  کی اس دائرہ کار میں مطلق زیادہ سے زیادہ اور مطلق کم سے کم قیمتیں پائی جائیں گی (شکل \حوالہ{شکل_کثیرالمتغیر_پہاڑی_وادی_الف} اور شکل \حوالہ{شکل_کثیرالمتغیر_پہاڑی_وادی_ب})۔ ان قیمتوں  کا جاننا اور ان نقطوں کا  جاننا،   جہاں یہ قیمتیں پائی جاتی ہیں،  ضروری ہے۔ہم جزوی تفرقات  سے  عموماً انہیں جان سکتے ہیں۔
./calculusMultivariableFunctionsAndPartialDerivatives.tex:6351:اگر تفاعل \عددی{f(x,y)} کے دائرہ کار کی اندرونی نقطہ \عددی{(a,b)} پر مقامی زیادہ سے زیادہ یا کم سے کم  قیمت پائی جاتی ہو، اور اگر تفاعل کے  یک رتبی تفرقات  موجود ہوں، تب \عددی{f_x(a,b)=0} اور \عددی{f_y(a,b)=0} ہوں گے۔ 
./calculusMultivariableFunctionsAndPartialDerivatives.tex:6437: پورا مستوی \عددی{xy} تفاعل \عددی{f} کا دائرہ کار  ہے ( لہٰذا کوئی سرحدی نقاط نہیں پائے جاتے ہیں) اور جزوی تفرقات \عددی{f_x=2x} اور \عددی{f_y=2y} ہر  نقطہ پر موجود ہوں گے۔ یوں مقامی انتہائی قیمتیں صرف اس صورت ممکن ہوں گی جب 
./calculusMultivariableFunctionsAndPartialDerivatives.tex:6473: پورا مستوی \عددی{xy} تفاعل \عددی{f} کا دائرہ کار   ہے ( لہٰذا کوئی سرحدی نقاط نہیں پائے جاتے ہیں) اور جزوی تفرقات \عددی{f_x=-2x} اور \عددی{f_y=2y} ہر  نقطہ پر موجود ہوں گے۔یوں مقامی انتہا قیمت صرف مبدا پر ممکن ہے۔ البتہ، مثبت \عددی{x} محور پر  \عددی{f} کی قیمت \عددی{f(x,0)=-x^2<0} اور مثبت \عددی{y} محور پر اس کی قیمت \عددی{f(0,y)=y^2>0} ہو گی۔ یوں  مستوی \عددی{xy} میں ہر  کھلا قرص، جس کا مرکز  \عددی{(0,0)} ہو، میں  ایسا نقاط پائے جاتے ہیں جہاں \عددی{f} مثبت ہو اور ایسے نقاط بھی پائے جاتے ہیں جہاں \عددی{f} منفی ہو۔  اس طرح مبدا پر مقامی انتہائی قیمت کے بجائے  تفاعل کا نقطہ زین پایا جاتا ہے (شکل \حوالہ{شکل_مثال_کثیرالمتغیر_نقطہ_زین_مبدا})۔ ہم دیکھتے ہیں کہ اس تفاعل کی  کوئی مقامی انتہائی قیمت نہیں پائی جاتی ہے۔
./calculusMultivariableFunctionsAndPartialDerivatives.tex:6477:ہم  دائرہ کار \عددی{R} کے اندرونی نقطہ \عددی{(a,b)} پر \عددی{f_x=f_y=0}  جانتے ہوئے یہ نہیں  کہہ سکتے ہیں کہ آیا اس نقطہ پر انتہائی قیمت پائی جاتی ہے۔ البتہ اگر \عددی{R} پر \عددی{f}  اور اس کے یک رتبی اور دو رتبی  جزوی تفرقات استمراری ہوں  تب ہم باقی معلومات درج ذیل مسئلہ (جسے حصہ \حوالہ{حصہ_کثیر_کلیہ_ٹیلر} میں ثابت کیا گیا ہے) سے دریافت کر سکتے ہیں۔
./calculusMultivariableFunctionsAndPartialDerivatives.tex:6481:فرض کریں کہ ایک قرص میں، جس  کا مرکز \عددی{(a,b)} ہے،  تفاعل \عددی{f} اور اس کے یک رتبی اور دو رتبی   جزوی تفرقات ہر نقطہ پر  استمراری ہیں  اور \عددی{f_x(a,b)=f_y(a,b)=0} ہے۔ اب
./calculusMultivariableFunctionsAndPartialDerivatives.tex:6649:  اگرچہ مسئلہ \حوالہ{مسئلہ_کثیرالمتغیر_انتہائی_قیمت_دو_رتبی_تفرقی_پرکھ}   طاقتور  ہے لیکن  اس کی کمزوری کو  نہ بھولیں۔یہ تفاعل کے دائرہ کار کے سرحدی نقطوں پر کارآمد نہیں ہوتا ہے جہاں تفاعل کے انتہائی قیمتیں اور غیر صفر تفرقات پائے جا  سکتے   ہیں۔  ساتھ ہی  اگر \عددی{f_x} یا \عددی{f_y} غیر موجود ہو تب بھی یہ مسئلہ کارآمد نہیں ہو گا۔
./calculusMultivariableFunctionsAndPartialDerivatives.tex:6659:اگر ایک قرص  میں، جس کا مرکز \عددی{(a,b)} ہو، ہر نقطہ پر \عددی{f}  کے یک رتبی اور دو رتبی جزوی تفرقات استمراری ہوں اور \عددی{f_x(a,b)=f_y(a,b)=0} ہو، آپ \عددی{f(a,b)}  کی  جماعت بندی \اصطلاح{دو رتبی تفرقی پرکھ} سے کر پائیں گے:
./calculusMultivariableFunctionsAndPartialDerivatives.tex:7042:(ا) کیا \عددی{f_x(a,b)=f_y(a,b)=0} ہوتے ہوئے ہر   صورت \عددی{(a,b)} پر \عددی{f} کا مقامی زیادہ سے  زیادہ قیمت  نقطہ یا کم سے کم قیمت  نقطہ پایا جائے گا؟ اپنے جواب کی وجہ پیش کریں۔   (ب)  اگر ایک قرص  میں ، جس کا مرکز  \عددی{(a,b)} ہو،  ہر نقطہ پر \عددی{f}  اور اس کے یک رتبی اور دو رتبی جزوی تفرقات استمراری ہوں، اور \عددی{f_{xx}(a,b)} اور \عددی{f_{yy}(a,b)} کی علامتیں ایک دوسرے سے مختلف ہوں  تب کیا \عددی{f} کے بارے میں کچھ کہنا ممکن ہو گا؟ اپنے جواب کی وجہ پیش کریں۔
./calculusMultivariableFunctionsAndPartialDerivatives.tex:7341:تفاعل کے دو رتبی تفرقات تلاش کر کے ممیز \عددی{f_{xx}f_{yy}-f_{xy}^2} معلوم کریں۔
./calculusMultivariableFunctionsAndPartialDerivatives.tex:8347: بشمول \عددی{\lambda_1} اور \عددی{\lambda_2} کے لحاظ سے، \عددی{h} کی تمام جزوی یک رتبی تفرقات حاصل کر کے \عددی{0} کے برابر رکھیں۔
./calculusMultivariableFunctionsAndPartialDerivatives.tex:8385:فرض کریں  خطہ \عددی{R} میں  تفاعل \عددی{f(x,y)}  کے  استمراری جزوی تفرقات  پائے جاتے ہیں اور اس خطہ میں نقطہ \عددی{N(a,b)}  پایا جاتا ہے جس پر \عددی{f_x=f_y=0} ہے۔ مزید فرض کریں کہ \عددی{h} اور \عددی{k} اتنے  چھوٹے ہیں کہ نقطہ \عددی{S(a+h,y+k)}اور \عددی{S}   سے \عددی{N} تک قطعے،   \عددی{R} میں پائے جاتے ہیں۔ہم قطعہ  \عددی{NS} کی مقدار معلوم مساوات لکھتے ہیں:
./calculusMultivariableFunctionsAndPartialDerivatives.tex:8394:چونکہ \عددی{f_x} اور \عددی{f_y}  قابل تفرق   ہیں (ان کے استمراری جزوی تفرقات پائے جاتے ہیں)     لہٰذا \عددی{t} کے لحاظ سے \عددی{F'} قابل تفرق ہو گا۔یوں درج ذیل لکھا جا سکتا ہے۔
./calculusMultivariableFunctionsAndPartialDerivatives.tex:8443:ہم فرض کرتے ہیں کہ بند مستطیل خطہ \عددی{R} جس کا مرکز \عددی{(x_0,y_0)}  ایک ایسے کھلے سلسلہ میں پایا جاتا ہے جس کے ہر نقطہ پر   \عددی{f}  کے استمراری دو رتبی جزوی تفرقات پائے جاتے ہیں۔ خطہ \عددی{R} میں \عددی{\abs{f_{xx}}}، \عددی{\abs{f_{yy}}} اور \عددی{\abs{f_{xy}}}  کی سب سے بڑی قیمت \عددی{B} ہے۔
./calculusMultivariableFunctionsAndPartialDerivatives.tex:8474:اگر ایک مستطیل، جس کا مرکز \عددی{(a,b)} ہو، میں  \عددی{f}  کے \عددی{n+1} رتبہ تک جزوی تفرقات ہر نقطہ پر  استمراری ہوں، تب ہم \عددی{F(t)} کے  کلیہ ٹیلر کو وسعت  دے کر
./calculusMultivariableFunctionsAndPartialDerivatives.tex:8482:حاصل کرتے ہیں۔ اس آخری  تسلسل   کے  دائیں ہاتھ    اولین \عددی{n} تفرقات میں   \عددی{t=0} پر  مساوات \حوالہ{مساوات_کثیرالمتغیر_عمومی_ٹیلر_عاملین} سے حاصل مطابقتی  تعلقات  پر   کر کے موزوں باقی جزو جمع کرتے ہوئے درج ذیل کلیہ حاصل ہو گا۔
./calculusMultivariableFunctionsAndPartialDerivatives.tex:8534:ہو گا۔ چونکہ تین رتبی تفرقات \عددی{\sin} اور \عددی{\cos} کے حاصل ضرب ہیں لہٰذا   ان کی  مطلق قیمتیں کسی صورت \عددی{1} سے زیادہ نہیں ہو سکتی ہیں۔ مزید \عددی{\abs{x}\le 0.1} اور \عددی{\abs{y}\le 0.1} ہیں لہٰذا
./calculusTranscendentalFunctions.tex:15:\حصہ{الٹ تفاعل اور ان کے تفرقات}\شناخت{حصہ_استعمال_تکمل_الٹ_تفاعل_اور_ان_کے_تفرق}
./calculusTranscendentalFunctions.tex:322:یہ تفرقات ایک دوسرے کے بالعکس متناسب ہیں۔ تفاعل \عددی{f} کی ترسیم لکیر \عددی{y=\tfrac{x}{2}+1} اور \عددی{f^{-1}} کی ترسیم لکیر \عددی{y=2x-2} ہے۔ ان لکیروں کے ڈھلوان ایک دوسرے کے بالعکس متناسب ہیں (شکل \حوالہ{شکل_ماورائی_عکس_میں_تفرق_تعلق})۔
./calculusTranscendentalFunctions.tex:2603:\موٹا{تفرقات}\\
./calculusVectorValuedFunctionsandMotionInSpace.tex:163:\جزوحصہء{تفرقات اور حرکت}
./calculusVectorValuedFunctionsandMotionInSpace.tex:278:چونکہ سمتی تفاعل کے تفرقات  جزو در جزو حاصل کرنا ممکن  ہے لہٰذا  سمتی تفاعل کے تفرقات کے قواعد کی نوعیت  غیر سمتی تفاعل کے تفرقات کے قواعد کی طرح ہو گی۔
./calculusVectorValuedFunctionsandMotionInSpace.tex:280:\موٹا{سمتی تفاعل کے تفرقات کے قواعد}\\
./calculusVectorValuedFunctionsandMotionInSpace.tex:328:ہو گا۔ ہم  شمار کنندہ  کے ساتھ   \عددی{\kvec{u}(t)\times \kvec{v}(t+h)} جمع اور منفی کرتے ہیں تا کہ درج بالا   کو ایسی حاصل تقسیمات  کی صورت میں لکھنا ممکن ہو جن میں  \عددی{\kvec{u}} اور \عددی{\kvec{v}} کے تفرقات    پائے جاتے ہوں۔یوں درج ذیل ہو گا۔
./calculusVectorValuedFunctionsandMotionInSpace.tex:409:اگر وقفہ \عددی{I}کے ہر نقطہ  پر \عددی{\tfrac{\dif \kvec{R}}{\dif t}=\kvec{r}} ہو تب    قابل تفرق سمتی تفاعل \عددی{\kvec{R}(t)}،وقفہ \عددی{I}   پر  سمتی تفاعل \عددی{\kvec{r}(t)} کا الٹ تفرق ہو گا۔ اگر \عددی{I} پر \عددی{\kvec{r}} کا الٹ تفرق \عددی{\kvec{R}} ہو تب ،ایک  وقت میں ایک جزو  کے ساتھ کام کرتے ہوئے،  یہ دکھایا جا سکتا ہے کہ \عددی{I} پر \عددی{\kvec{r}} کے الٹ تفرق  کی صورت \عددی{\kvec{R}+\kvec{C}} ہو گی جہاں \عددی{\kvec{C}} کوئی مستقل سمتیہ ہو گا (سوال \حوالہ{سوال_سمتی_تفاعل_الٹ_تفرقات})۔وقفہ \عددی{I} پر \عددی{\kvec{r}} کے الٹ تفرقات کا سلسلہ \عددی{I} پر \عددی{\kvec{r}} کا \اصطلاح{غیر قطعی تکمل}\فرہنگ{تکمل!غیر قطعی}\حاشیہب{indefinite integral}\فرہنگ{integral!indefinite} ہو گا۔ 
./calculusVectorValuedFunctionsandMotionInSpace.tex:412:متغیر \عددی{t} کے لحاظ سے \عددی{\kvec{r}} کا غیر قطعی تکمل، \عددی{\kvec{r}} کے تمام الٹ تفرقات کا سلسلہ ہو گا، جس کو \عددی{\int\kvec{r}(t)\dif t} سے ظاہر کیا جاتا ہے۔ اگر \عددی{\kvec{r}} کا الٹ تفرق \عددی{\kvec{R}} ہو تب درج ذیل ہو گا۔
./calculusVectorValuedFunctionsandMotionInSpace.tex:1080:\ابتدا{سوال}\شناخت{سوال_سمتی_تفاعل_الٹ_تفرقات}\ترچھا{سمتی تفاعل کے الٹ تفرقات}\\
./calculusVectorValuedFunctionsandMotionInSpace.tex:1083:غیر سمتی تفاعل کے مسئلہ اوسط قیمت کا ضمنی نتیجہ  \عددی{2} استعمال کرتے ہوئے  دکھائیں کہ اگر  وقفہ \عددی{I} پر دو سمتی تفاعل \عددی{\kvec{R}_1(t)} اور \عددی{\kvec{R}_2(t)}  کے تفرقات  متماثل  ہوں  تب پورے  \عددی{I} پر ان تفاعل میں صرف ایک مستقل  سمتی قیمت کا فرق ہو گا۔
./calculusVectorValuedFunctionsandMotionInSpace.tex:1796:چونکہ مساوات \حوالہ{مساوات_سمتی_تفاعل_لمبائی_قوس_ت} میں جذر کے اندر تفرقات  استمراری   (ہموار منحنی)  ہیں احصاء کے بنیادی مسئلہ کے تحت \عددی{s} متغیر \عددی{t} کا قابل تفرق تفاعل ہو گا     اور یہ تفرق درج ذیل ہو گا۔
./calculusVectorValuedFunctionsandMotionInSpace.tex:2107:اس حصہ میں ہم   تین آپس میں عمودی اکائی سمتیات پر مبنی ایسا چوکھٹ متعارف کرتے ہیں جو  فضا میں منحنی پر جسم کے ساتھ ساتھ چلتا ہو (شکل \حوالہ{شکل_سمتی_تفاعل_چھوکٹ}) ۔ اس چوکھٹ کے تین سمتیات  ہیں۔ پہلا  اکائی مماسی سمتیہ \عددی{\kvec{T}} ہے۔ دوسرا \عددی{\kvec{N}} ہے جو  \عددی{\tfrac{\dif \kvec{R}}{\dif s}} کے رخ اکائی سمتیہ ہے۔تیسرا اکائی سمتیہ  \عددی{\kvec{B}=\kvec{T}\times\kvec{N}} ہے ۔ یہ سمتیات اور ان کے تفرقات  اگر معلوم ہوں، فضا  میں سواری  کی  سمت بندی اور اس کی راہ  میں موڑ  اور بل  کے بارے میں مفید معلومات   مہیا کرتے ہیں۔
./calculusVectorValuedFunctionsandMotionInSpace.tex:3772:فرض کریں  ایک مستوی میں متحرک ذرے کا تعین گر سمتیہ \عددی{\kvec{r}} ہے اور  یہ  سمتیہ \عددی{\tfrac{\dif S}{\dif t}} کی شرح سے رقبہ واضح کرتا ہے۔محدد متعارف کئے بغیر اور مطلوبہ تفرقات  کی موجودگی تصور کرتے ہوئے، بڑھوتری اور حد  پر مبنی  درج ذیل مساوات کی  جیومیٹریائی جواز پیش کریں۔ 
./calculusVectorValuedFunctionsandMotionInSpace.tex:3828:یہ فرض کرتے ہوئے  کہ \عددی{t} کے لحاظ سے درکار تفرقات موجود ہیں، \عددی{\kvec{v}=\kvec{r}} اور \عددی{\kvec{a}=\ddot{\kvec{r}}} کو \عددی{\kvec{u}_r}، \عددی{\kvec{u}_{\theta}}، \عددی{\ak}، \عددی{\dot{r}} اور \عددی{\dot{\theta}} کی صورت میں لکھیں۔(نقطہ \عددی{t} کے لحاظ سے تفرق کو ظاہر کرتا ہے، لہٰذا \عددی{\dot{\kvec{r}}} سے مراد \عددی{\tfrac{\dif\kvec{r}}{\dif t}} اور \عددی{\ddot{\kvec{r}}} سے مراد \عددی{\tfrac{\dif^{\,2}\kvec{r}}{\dif t^2}} ہو گا، وغیرہ وغیرہ۔)  حصہ \حوالہ{حصہ_سمتی_تفاعل_فلکی_سیاروں_اور_مصنوعی_سیاروں_کی_حرکت} میں ان کلیات کو اخذ کیا گیا ہے  اور ساتھ ہی ان سمتیات کو فلکیاتی سیاروں کی حرکت بیان کرنے کے لئے استعمال کیا گیا ہے۔
